% 한글 레이텍스 두단 논문 초안
\documentclass[11pt, twocolumn]{article}
\usepackage[utf-8]{inputenc}
\usepackage[hangul]{babel}
\usepackage{kotex}
\usepackage{geometry}
\usepackage{hyperref}
\usepackage{graphicx}
\usepackage{amsmath}
\usepackage{amssymb}
\usepackage{cite}
\usepackage{xcolor}

% 페이지 설정
\geometry{
    top=2cm,
    bottom=2cm,
    left=1.5cm,
    right=1.5cm,
    columnsep=0.8cm
}

% 제목 스타일
\title{의미론적 숲 모델에 대한 리뷰 대응: \\
해석 가능성, 하이브리드 아키텍처, 성능 분해 분석}
\author{ACO Semantic Forest Team}
\date{\today}

\begin{document}

\maketitle

\section{개요}

본 문서는 의미론적 숲(Semantic Forest) 모델의 핵심 설계 선택과 성능 분석에 대한 리뷰어의 예상 질문에 대한 답변 초안입니다. 
네 가지 주요 질문에 대해 이론적 근거, 실험 증거, 그리고 추가 분석 방향을 제시합니다.

\section{Q1: Neural Network 기반 Tree를 사용하지 않은 이유}

\subsection{핵심 입장}

본 연구의 주요 목표는 예측 성능 극대화가 아닌 \textbf{해석 가능성과 온톨로지 정합성}에 있습니다. 
Neural-tree 계열 모델(예: SAINT, XGBoost with neural splits)은 높은 성능을 제공하지만, 다음의 단점이 있습니다:

\begin{itemize}
    \item 모델 복잡도 증가에 따른 규칙 추적 어려움
    \item 노드 단위 의미 해석 불명확
    \item 의료/과학 도메인에서의 신뢰성 감소
\end{itemize}

\subsection{본 방법의 장점}

반면 Semantic-rule tree 기반 접근은:

\begin{enumerate}
    \item \textbf{규칙의 투명성}: 각 분할이 온톨로지 기반 개념과 직접 연결
    \item \textbf{제한 데이터 구간 안정성}: 중소 규모 데이터셋에서 재현성 우수
    \item \textbf{계산 효율성}: 재현 가능한 학습, 모듈식 확장 가능
\end{enumerate}

\subsection{제안된 검증 방법}

\begin{itemize}
    \item 복잡도-성능-해석성 3축 트레이드오프 표
    \item 다중 랜덤 시드 실험으로 안정성 검증
    \item 규칙 충실도(faithfulness) 사례 분석
\end{itemize}

\section{Q2: 딥러닝 모델과의 결합 가능성 및 시각화 차이}

\subsection{하이브리드 아키텍처 제안}

성능 개선을 위해 두 단계 파이프라인을 제시합니다:

\[
\text{입력} \xrightarrow{\text{Deep Encoder}} \text{표현} \xrightarrow{\text{Semantic Tree}} \text{규칙 경로} \rightarrow \text{의사결정}
\]

깊은 신경망이 복잡한 패턴을 압축하고, 트리가 최종 판단을 투명한 규칙으로 제시합니다.

\subsection{시각화 비교}

\begin{table}[h]
    \centering
    \small
    \begin{tabular}{|l|l|l|}
    \hline
    \textbf{기준} & \textbf{딥러닝} & \textbf{Semantic Tree / 하이브리드} \\
    \hline
    중심 설명 & saliency, attribution & 규칙 경로, 개념 기여도 \\
    \hline
    입도(granularity) & 픽셀/특징 & 의사결정 노드/규칙 \\
    \hline
    도메인 정합성 & 낮음 & 높음 \\
    \hline
    \end{tabular}
\end{table}

\subsection{기대 효과}

하이브리드 모델에서는 두 가지 설명 채널을 동시 제시:
\begin{enumerate}
    \item 인코더의 특징 중요도 (표현 수준)
    \item 트리의 규칙 경로 (개념 수준)
\end{enumerate}

\section{Q3: 대규모 데이터에서 성능 저하 시 설명 가능성}

\subsection{성능 저하가 항상 부정적인가?}

아니오. 성능 저하가 발생해도 설명 시각화는 오히려 더 중요해집니다.

\subsection{원인 분해 분석}

성능 저하를 단순 알고리즘 한계가 아닌 \textbf{구조적 원인}으로 분석:

\begin{itemize}
    \item \textbf{분포 이동(Distribution shift)}: 데이터 특성 변화 추적
    \item \textbf{레이블 노이즈}: 신뢰도 낮은 샘플 식별
    \item \textbf{의미론 커버리지}: 온톨로지 개념 활용도 분석
    \item \textbf{모델 용량 부족}: 디이터 복잡도 vs. 모델 복잡도 비교
\end{itemize}

\subsection{제안된 시각화}

\begin{enumerate}
    \item 학습 곡선: 데이터 크기별 성능 추세
    \item 오류 분해: 클래스/부분군별 실패 패턴
    \item 개념 맵: 온톨로지 영역별 규칙 밀도
    \item 신뢰도 대역: 예측 불확실성 시각화
\end{enumerate}

\section{Q4: 규칙 분할 기준과 매개변수 기여도}

\subsection{단순 보팅인가?}

절대 아닙니다. 분할 결정은 \textbf{다항 목적함수} 기반입니다:

\[
\text{Objective} = \alpha \cdot \text{Gain} + \beta \cdot \text{Semantic\_Fit} + \gamma \cdot \text{Regularization}
\]

여기서:
\begin{itemize}
    \item $\text{Gain}$: 정보 이득, 불순도 감소
    \item $\text{Semantic\_Fit}$: 온톨로지 정합성
    \item $\text{Regularization}$: 복잡도 페널티
\end{itemize}

\subsection{규칙 기여도 정량화}

각 규칙의 영향력을 네 가지 방식으로 측정:

\begin{enumerate}
    \item \textbf{경로 누적 기여도}: 해당 규칙 포함 경로의 예측 개선량
    \item \textbf{부트스트랩 안정성}: 반복 학습에서 선택 빈도
    \item \textbf{제거 영향도}: 규칙 제거 시 성능 저하 폭
    \item \textbf{특징 중요도 변화}: SHAP, permutation importance 등
\end{enumerate}

\subsection{결론}

규칙은 정량적 다목적 최적화 기반으로 선택되며, 단순 성능 보팅이 아닌 
\textbf{예측력 + 의미론적 타당성}의 균형 적용입니다.

\section{향후 실험 계획}

\begin{enumerate}
    \item 하이브리드 구조 구현 및 성능 비교 (DL pure vs. 본 방법 vs. 하이브리드)
    \item 대규모 데이터 성능 분해 분석
    \item 규칙 기여도 대시보드 구현
    \item 도메인 전문가 설명력 평가 (정성 실험)
\end{enumerate}

\section{결론}

본 모델의 설계는 최고 성능보다 \textbf{해석 가능성과 신뢰성 중심} 선택입니다. 
각 리뷰 질문에 대해 구체적 이론적 근거, 예상 실험 증거, 그리고 성능-설명력 균형 전략을 제시했습니다.

\end{document}
